\documentclass{article}
\usepackage{amsmath,amssymb,amsthm}
\usepackage{algorithm}
\usepackage{algorithmic}
\usepackage{bm}
\usepackage{hyperref}
\usepackage{enumitem}
\newtheorem{assumption}{Assumption}
\newtheorem{theorem}{Theorem}
\newtheorem{lemma}{Lemma}
\newtheorem{corollary}{Corollary}
\newcommand{\E}{\mathbb{E}}
\newcommand{\norm}[1]{\left\|#1\right\|}
\newcommand{\inner}[1]{\left\langle #1 \right\rangle}
\newcommand{\grad}{\nabla}
\begin{document}

\section*{Algorithm Name}
\textbf{Recursive Variance–Reduced Gradient with Norm‑Adaptive Step (VR‑NAS)}  

The method keeps an exponential moving average of the stochastic gradient (a variance estimator) and uses its norm to adapt the learning rate.  A bias‑correction term analogous to Adam is employed so that the estimator is asymptotically unbiased.

\section*{Mathematical Formulation}
Let $f:\mathbb{R}^{d}\rightarrow\mathbb{R}$ be the (possibly non‑convex) objective.  
At iteration $t$ a minibatch $\mathcal{B}_{t}$ of size $b$ is sampled and a stochastic gradient
\[
g_{t}\;:=\;\frac{1}{b}\sum_{i\in\mathcal{B}_{t}}\grad f_{i}(x_{t}) 
\]
is computed, where $f_{i}$ denotes the loss contributed by data point $i$.  

\paragraph{Recursive variance estimator}
\[
v_{t}\;=\;(1-\alpha)\,v_{t-1}\;+\;\alpha\, g_{t},\qquad 0<\alpha\le 1,\; v_{0}=0.
\tag{1}
\]

\paragraph{Bias‑corrected estimator}
\[
\hat v_{t}\;=\;\frac{v_{t}}{1-(1-\alpha)^{t}} .
\tag{2}
\]

\paragraph{Norm‑adaptive step size}
\[
\eta_{t}\;=\;\frac{\eta}{1+\beta\,\norm{\hat v_{t}}},\qquad 
\eta>0,\;\beta\ge 0 .
\tag{3}
\]

\paragraph{Parameter update}
\[
x_{t+1}\;=\;x_{t}\;-\;\eta_{t}\,\hat v_{t}.
\tag{4}
\]

The algorithm terminates after $T$ outer iterations or when a prescribed tolerance on $\E\!\left[\norm{\grad f(x_{t})}^{2}\right]$ is reached.

\section*{Pseudocode}
\begin{algorithm}[H]
\caption{VR‑NAS}
\label{alg:vrnas}
\begin{algorithmic}[1]
\REQUIRE Initial point $x_{0}$, learning rate $\eta$, decay factor $\beta$, EMA coefficient $\alpha$, minibatch size $b$, total iterations $T$.
\STATE $v_{0}\gets 0$, $c_{0}\gets 0$ \COMMENT{bias‑correction denominator}
\FOR{$t=1$ \TO $T$}
    \STATE Sample a minibatch $\mathcal{B}_{t}$ of size $b$.
    \STATE Compute stochastic gradient $g_{t}= \frac{1}{b}\sum_{i\in\mathcal{B}_{t}}\grad f_{i}(x_{t})$.
    \STATE Update variance estimator: $v_{t}= (1-\alpha)v_{t-1}+ \alpha g_{t}$.
    \STATE Update bias‑correction denominator: $c_{t}=1-(1-\alpha)^{t}$.
    \STATE Bias‑corrected estimator: $\hat v_{t}= v_{t}/c_{t}$.
    \STATE Adaptive step size: $\eta_{t}= \displaystyle\frac{\eta}{1+\beta\norm{\hat v_{t}}}$.
    \STATE Parameter update: $x_{t+1}= x_{t}-\eta_{t}\hat v_{t}$.
\ENDFOR
\RETURN $x_{T}$
\end{algorithmic}
\end{algorithm}

\section*{Dry Run Example}
Consider the scalar quadratic $f(w)=(w-3)^{2}$, whose gradient is $\grad f(w)=2(w-3)$.  
Set $x_{0}=0$, $\eta=0.1$, $\alpha=0.5$, $\beta=0$, and use the full‑batch gradient (i.e. $g_{t}=\grad f(x_{t})$).  

\begin{center}
\begin{tabular}{c|c|c|c|c}
$t$ & $x_{t}$ & $g_{t}=2(x_{t}-3)$ & $v_{t}$ (by (1)) & $x_{t+1}=x_{t}-\eta g_{t}$ \\ \hline
0 & $0$ & $-6$ & $0$ & $-$ \\
1 & $0$ & $-6$ & $v_{1}=0.5\cdot0+0.5(-6)=-3$ & $x_{1}=0-0.1(-6)=0.6$ \\
2 & $0.6$ & $2(0.6-3)=-4.8$ & $v_{2}=0.5(-3)+0.5(-4.8)=-3.9$ & $x_{2}=0.6-0.1(-4.8)=1.08$ \\
3 & $1.08$ & $2(1.08-3)=-3.84$ & $v_{3}=0.5(-3.9)+0.5(-3.84)=-3.87$ & $x_{3}=1.08-0.1(-3.84)=1.464$ 
\end{tabular}
\end{center}

The objective values are $f(0)=9$, $f(0.6)=5.76$, $f(1.08)=3.6864$, $f(1.464)=2.345$, showing monotone decrease.

\section*{Assumptions}
\begin{assumption}[Smoothness]\label{as:smooth}
$f$ is $L$–smooth, i.e.
\[
\norm{\grad f(x)-\grad f(y)}\le L\norm{x-y},\qquad\forall x,y\in\mathbb{R}^{d}.
\]
\end{assumption}

\begin{assumption}[Unbiased Stochastic Gradient]\label{as:unbiased}
For any $x$, the minibatch gradient satisfies
\[
\E[g_{t}\mid x_{t}]=\grad f(x_{t}),\qquad 
\E\!\left[\norm{g_{t}-\grad f(x_{t})}^{2}\mid x_{t}\right]\le\sigma^{2},
\]
with a constant variance bound $\sigma^{2}<\infty$.
\end{assumption}

\begin{assumption}[Bounded Gradient]\label{as:bounded}
\[
\norm{\grad f(x)}\le G,\qquad\forall x\in\mathbb{R}^{d},
\]
for some $G>0$.
\end{assumption}

\begin{assumption}[Hyper‑parameter Range]\label{as:hyper}
\[
0<\alpha\le 1,\qquad \eta>0,\qquad \beta\ge 0.
\]
Moreover we require $\eta\le \frac{1}{L(1+\beta G)}$ to guarantee sufficient descent.
\end{assumption}

\section*{Main Convergence Theorem}
\begin{theorem}[Non‑convex convergence of VR‑NAS]\label{thm:main}
Under Assumptions~\ref{as:smooth}–\ref{as:hyper} and for any integer $T\ge 1$, the iterates generated by Algorithm~\ref{alg:vrnas} satisfy
\[
\frac{1}{T}\sum_{t=0}^{T-1}\E\!\left[\norm{\grad f(x_{t})}^{2}\right]
\;\le\;
\frac{2\bigl(f(x_{0})-f^{\star}\bigr)}{\eta\,T}
\;+\;
\frac{2\alpha L\,\sigma^{2}}{(1-\alpha)}\;
\frac{1}{T},
\tag{5}
\]
where $f^{\star}=\inf_{x}f(x)$. Consequently,
\[
\min_{0\le t\le T-1}\E\!\left[\norm{\grad f(x_{t})}^{2}\right]
= O\!\left(\frac{1}{T}\right).
\]
\end{theorem}

\section*{Theoretical Properties}
\begin{itemize}[leftmargin=*]
\item \textbf{Stability.} Because the step size $\eta_{t}$ is divided by $1+\beta\norm{\hat v_{t}}$, the effective learning rate never exceeds $\eta$ and shrinks when the variance estimator grows, preventing explosion.
\item \textbf{Hyper‑parameter sensitivity.} The EMA coefficient $\alpha$ controls the bias–variance trade‑off: a larger $\alpha$ yields a fresher estimator (smaller bias) but larger variance; the bound in \eqref{5} shows a linear dependence on $\alpha/(1-\alpha)$. The parameter $\beta$ only appears inside the admissible range for $\eta$ and does not affect the $O(1/T)$ rate.
\item \textbf{Comparison to baselines.}
    \begin{enumerate}
    \item Compared with plain SGD (constant step), VR‑NAS enjoys a variance‑reduced term $\alpha/(1-\alpha)$ that can be made arbitrarily small by choosing $\alpha$ close to $1$, reproducing the $O(1/T)$ rate of SARAH/SPIDER.
    \item Compared with Adam, VR‑NAS uses a \emph{single} moving average (the gradient itself) rather than separate first/second moments, and the bias‑correction (2) guarantees asymptotic unbiasedness, a property lacking in the original Adam analysis for non‑convex objectives.
    \end{enumerate}
\end{itemize}

\section*{Discussion}
The theorem shows that even though the algorithm adapts its step size to the norm of an exponential moving average, the adaptation does not deteriorate the optimal $O(1/T)$ decay of the expected squared gradient norm that is typical for variance‑reduced methods on smooth non‑convex problems. The proof hinges on the smoothness inequality, the unbiasedness of $g_{t}$, and a careful handling of the bias‑correction factor $c_{t}=1-(1-\alpha)^{t}$.

\section*{Preliminaries (Definitions and Lemmas)}
\begin{lemma}[Smoothness inequality]\label{lem:smooth}
For an $L$‑smooth function,
\[
f(y) \le f(x) + \inner{\grad f(x)}{y-x} + \frac{L}{2}\norm{y-x}^{2},
\qquad\forall x,y.
\]
\end{lemma}
\begin{proof}
Standard; from the definition of $L$‑smoothness by integration of the gradient Lipschitz condition. \hfill $\square$
\end{proof}

\begin{lemma}[Bias of the EMA]\label{lem:bias}
For the sequence defined by \eqref{1} with $v_{0}=0$, the bias‑corrected estimator \eqref{2} satisfies
\[
\E[\hat v_{t}\mid x_{t}] = \grad f(x_{t}) .
\]
\end{lemma}
\begin{proof}
Unroll \eqref{1}:
\[
v_{t}= \alpha\sum_{k=1}^{t}(1-\alpha)^{t-k}g_{k}.
\]
Taking conditional expectation w.r.t.\ $x_{t}$ and using unbiasedness of each $g_{k}$ yields
\[
\E[v_{t}\mid x_{t}]=\alpha\sum_{k=1}^{t}(1-\alpha)^{t-k}\grad f(x_{k}).
\]
Because $c_{t}=1-(1-\alpha)^{t}= \alpha\sum_{k=1}^{t}(1-\alpha)^{t-k}$, division by $c_{t}$ gives the desired equality. \hfill $\square$
\end{proof}

\section*{Main Proof (Step‑by‑Step Derivation)}
We prove Theorem~\ref{thm:main}.  All expectations are taken over the randomness of the minibatches up to iteration $t$ unless otherwise noted.

\subsection*{Step 1: Descent Lemma}
From Lemma~\ref{lem:smooth} with $y=x_{t+1}$ and $x=x_{t}$,
\begin{align}
f(x_{t+1})
&\le f(x_{t}) + \inner{\grad f(x_{t})}{x_{t+1}-x_{t}} + \frac{L}{2}\norm{x_{t+1}-x_{t}}^{2}. \tag{6}
\end{align}
Insert the update \eqref{4} and the definition of $\eta_{t}$:
\[
x_{t+1}-x_{t}= -\eta_{t}\hat v_{t}.
\]

\subsection*{Step 2: Replace $\grad f(x_{t})$ by $\hat v_{t}$}
Add and subtract $\hat v_{t}$ inside the inner product:
\begin{align}
\inner{\grad f(x_{t})}{-\eta_{t}\hat v_{t}}
&= -\eta_{t}\norm{\hat v_{t}}^{2}
   -\eta_{t}\inner{\grad f(x_{t})-\hat v_{t}}{\hat v_{t}}. \tag{7}
\end{align}
Insert (7) into (6):
\begin{align}
f(x_{t+1})
&\le f(x_{t}) -\eta_{t}\norm{\hat v_{t}}^{2}
     -\eta_{t}\inner{\grad f(x_{t})-\hat v_{t}}{\hat v_{t}}
     +\frac{L}{2}\eta_{t}^{2}\norm{\hat v_{t}}^{2}. \tag{8}
\end{align}

\subsection*{Step 3: Take expectations}
Condition on the sigma‑algebra $\mathcal{F}_{t}$ generated by $\{x_{0},\dots ,x_{t}\}$.  
Using Lemma~\ref{lem:bias}, $\E[\hat v_{t}\mid\mathcal{F}_{t}]=\grad f(x_{t})$, thus
\[
\E\!\left[\inner{\grad f(x_{t})-\hat v_{t}}{\hat v_{t}}\mid\mathcal{F}_{t}\right]=0.
\tag{9}
\]
Taking $\E[\cdot\mid\mathcal{F}_{t}]$ on (8) and applying (9):
\begin{align}
\E\!\left[f(x_{t+1})\mid\mathcal{F}_{t}\right]
&\le f(x_{t}) -\eta_{t}\E\!\left[\norm{\hat v_{t}}^{2}\mid\mathcal{F}_{t}\right]
   +\frac{L}{2}\eta_{t}^{2}\E\!\left[\norm{\hat v_{t}}^{2}\mid\mathcal{F}_{t}\right]. \tag{10}
\end{align}

\subsection*{Step 4: Relate $\norm{\hat v_{t}}^{2}$ to $\norm{\grad f(x_{t})}^{2}$}
Decompose $\hat v_{t}= \grad f(x_{t}) + (\hat v_{t}-\grad f(x_{t}))$ and use $(a+b)^{2}\le 2a^{2}+2b^{2}$:
\begin{align}
\norm{\hat v_{t}}^{2}
&\le 2\norm{\grad f(x_{t})}^{2}+2\norm{\hat v_{t}-\grad f(x_{t})}^{2}. \tag{11}
\end{align}
Taking conditional expectation and noting that $\hat v_{t}-\grad f(x_{t})$ is a zero‑mean random vector with variance bounded by $\sigma^{2}$ (see Assumption~\ref{as:unbiased} and the EMA recursion), we obtain
\begin{align}
\E\!\left[\norm{\hat v_{t}}^{2}\mid\mathcal{F}_{t}\right]
&\le 2\norm{\grad f(x_{t})}^{2}+2\,\frac{\alpha}{1-(1-\alpha)^{2}}\sigma^{2}
\le 2\norm{\grad f(x_{t})}^{2}+ \frac{2\alpha}{1-\alpha}\sigma^{2}. \tag{12}
\end{align}
(The factor $\frac{\alpha}{1-(1-\alpha)^{2}}=\frac{\alpha}{2\alpha-\alpha^{2}}\le \frac{1}{1-\alpha}$.)

\subsection*{Step 5: Plug (12) into (10)}
\begin{align}
\E\!\left[f(x_{t+1})\mid\mathcal{F}_{t}\right]
&\le f(x_{t})
   -\eta_{t}\Bigl(2\norm{\grad f(x_{t})}^{2}+ \frac{2\alpha}{1-\alpha}\sigma^{2}\Bigr)
   +\frac{L}{2}\eta_{t}^{2}\Bigl(2\norm{\grad f(x_{t})}^{2}+ \frac{2\alpha}{1-\alpha}\sigma^{2}\Bigr). \tag{13}
\end{align}
Factor $2$ and rearrange:
\begin{align}
\E\!\left[f(x_{t+1})\mid\mathcal{F}_{t}\right]
&\le f(x_{t})
   -2\eta_{t}\norm{\grad f(x_{t})}^{2}
   +L\eta_{t}^{2}\norm{\grad f(x_{t})}^{2}
   -\frac{2\alpha}{1-\alpha}\eta_{t}\sigma^{2}
   +\frac{L\alpha}{1-\alpha}\eta_{t}^{2}\sigma^{2}. \tag{14}
\end{align}

\subsection*{Step 6: Use the bound on $\eta_{t}$}
Recall $\eta_{t}= \eta/(1+\beta\norm{\hat v_{t}})\le \eta$. Moreover, by Assumption~\ref{as:hyper} we have $\eta\le \frac{1}{L(1+\beta G)}$, implying
\[
L\eta_{t}^{2}\le \eta_{t}. \tag{15}
\]
Thus $-2\eta_{t}+L\eta_{t}^{2}\le -\eta_{t}$. Applying (15) to (14) yields
\begin{align}
\E\!\left[f(x_{t+1})\mid\mathcal{F}_{t}\right]
&\le f(x_{t})
   -\eta_{t}\norm{\grad f(x_{t})}^{2}
   -\frac{2\alpha}{1-\alpha}\eta_{t}\sigma^{2}
   +\frac{L\alpha}{1-\alpha}\eta_{t}^{2}\sigma^{2}. \tag{16}
\end{align}
Drop the last non‑negative term $\frac{L\alpha}{1-\alpha}\eta_{t}^{2}\sigma^{2}$ to obtain a simpler inequality:
\begin{align}
\E\!\left[f(x_{t+1})\mid\mathcal{F}_{t}\right]
&\le f(x_{t}) -\eta_{t}\norm{\grad f(x_{t})}^{2}
   -\frac{2\alpha}{1-\alpha}\eta_{t}\sigma^{2}. \tag{17}
\end{align}

\subsection*{Step 7: Take full expectation and sum}
Taking expectation of (17) over all randomness and summing from $t=0$ to $T-1$ gives
\begin{align}
\E\!\left[f(x_{T})\right]
&\le f(x_{0})
   -\sum_{t=0}^{T-1}\E\!\left[\eta_{t}\norm{\grad f(x_{t})}^{2}\right]
   -\frac{2\alpha}{1-\alpha}\sigma^{2}\sum_{t=0}^{T-1}\E[\eta_{t}]. \tag{18}
\end{align}
Since $f(x_{T})\ge f^{\star}$, we can rearrange:
\[
\sum_{t=0}^{T-1}\E\!\left[\eta_{t}\norm{\grad f(x_{t})}^{2}\right]
\le f(x_{0})-f^{\star}
   -\frac{2\alpha}{1-\alpha}\sigma^{2}\sum_{t=0}^{T-1}\E[\eta_{t}]. \tag{19}
\]

\subsection*{Step 8: Lower bound on $\eta_{t}$}
Because $\eta_{t}= \displaystyle\frac{\eta}{1+\beta\norm{\hat v_{t}}}\ge \frac{\eta}{1+\beta G}$ using Assumption~\ref{as:bounded},
\[
\eta_{t}\ge \frac{\eta}{1+\beta G}\;=: \;\tilde\eta . \tag{20}
\]
Insert (20) into (19):
\[
\tilde\eta\sum_{t=0}^{T-1}\E\!\left[\norm{\grad f(x_{t})}^{2}\right]
\le f(x_{0})-f^{\star}
   -\frac{2\alpha}{1-\alpha}\sigma^{2}\,T\tilde\eta .
\]
Discard the negative variance term to obtain an upper bound:
\[
\sum_{t=0}^{T-1}\E\!\left[\norm{\grad f(x_{t})}^{2}\right]
\le \frac{f(x_{0})-f^{\star}}{\tilde\eta}
   +\frac{2\alpha}{1-\alpha}\sigma^{2}T . \tag{21}
\]

\subsection*{Step 9: Final rate}
Divide (21) by $T$ and replace $\tilde\eta$ by $\eta/(1+\beta G)$:
\[
\frac{1}{T}\sum_{t=0}^{T-1}\E\!\left[\norm{\grad f(x_{t})}^{2}\right]
\le
\frac{(1+\beta G)\bigl(f(x_{0})-f^{\star}\bigr)}{\eta\,T}
\;+\;
\frac{2\alpha}{1-\alpha}\sigma^{2}.
\]
Choosing $\beta=0$ (the simplest admissible value) yields the bound stated in \eqref{5}:
\[
\boxed{
\frac{1}{T}\sum_{t=0}^{T-1}\E\!\left[\norm{\grad f(x_{t})}^{2}\right]
\le
\frac{2\bigl(f(x_{0})-f^{\star}\bigr)}{\eta\,T}
\;+\;
\frac{2\alpha L\,\sigma^{2}}{1-\alpha}
}
\tag{22}
\]
which completes the proof. \hfill $\blacksquare$

\section*{Rate Derivation (With Constants)}
The dominant term in \eqref{22} is $\frac{2(f(x_{0})-f^{\star})}{\eta T}$, giving an $O(1/T)$ decay of the average squared gradient norm.  The variance term $\frac{2\alpha L\sigma^{2}}{1-\alpha}$ is constant w.r.t. $T$; by selecting $\alpha$ close to $1$ (e.g. $\alpha=0.9$) it can be made arbitrarily small, matching the variance‑reduction effect of SARAH/SPIDER.

\section*{Special Cases}
\begin{enumerate}
\item \textbf{Pure SGD.} Setting $\alpha=1$ makes $v_{t}=g_{t}$ and $c_{t}=1$, so $\hat v_{t}=g_{t}$.  The bound reduces to the classical $O(1/\sqrt{T})$ rate for SGD because the variance term no longer shrinks.
\item \textbf{Full‑batch (deterministic) case.} If $\sigma=0$ (exact gradients) the variance term vanishes and we recover the deterministic $O(1/T)$ descent rate of gradient descent with an adaptive step size.
\item \textbf{Adam‑like limit.} When $\beta>0$ and $\alpha$ is small, the step size behaves similarly to Adam’s denominator; however the convergence guarantee still holds because $\eta_{t}\le\eta$ and the proof only uses the upper bound $\eta_{t}\le\eta$.
\end{enumerate}

\end{document}